\chapter{函数极限与实数基本定理}

\newpage

\section{求极限值}

\newpage

\subsection{等价无穷小及已知极限}

\begin{proposition}[\marktwostar\ 常用的等价关系与极限]
当 $x\to0$ 时,
\[
    \begin{aligned}
        x&\sim\sin x\sim\tan x\sim\arcsin x\sim\arctan x\sim\ln(1+x)\\
        &\sim e^x-1
        \sim\dfrac{a^x-1}{\ln a}
        \sim\dfrac{(1+x)^b-1}{b},
        \qquad a>0,\ b\neq0.
    \end{aligned}
\]
此外,
\[
    1-\cos x\sim\dfrac{1}{2}x^2,
    \qquad
    (1+x)^{\frac{1}{x}}
    =e-\dfrac{e}{2}x+\dfrac{11e}{24}x^2+o(x^2).
\]
又有
\[
    1+\dfrac{1}{2}+\dfrac{1}{3}+\cdots+\dfrac{1}{n}
    =\ln n+\gamma+\alpha_n,
    \qquad
    \lim\limits_{n\to\infty}\alpha_n=0,
\]
其中 $\gamma$ 为 Euler 常数.并且
\[
    n!=\sqrt{2\pi n}\left(\dfrac{n}{e}\right)^n
    e^{\frac{\theta_n}{12n}}
    \sim
    \sqrt{2\pi n}\left(\dfrac{n}{e}\right)^n,
    \qquad 0\leq\theta_n\leq1.
\]
\end{proposition}

\begin{example}
求下列极限:
\begin{enumerate}[label=(\arabic*)]
    \item $\displaystyle
    \lim\limits_{n\to\infty}
    \dfrac{n^3\sqrt[n]{2\left(1-\cos\dfrac{1}{n^2}\right)}}{\sqrt{n^2+1}-n}$;
    \item $\displaystyle
    \lim\limits_{x\to0}
    \dfrac{x(1-\cos x)}{(1-e^x)\sin x^2}$;
    \item $\displaystyle
    \lim\limits_{x\to0}
    \dfrac{\arctan x}{\ln(1+\sin x)}$.
\end{enumerate}
\end{example}
\exampleblank

\begin{example}
求
\[
    \lim\limits_{n\to\infty}
    \left(\dfrac{1}{n+1}+\dfrac{1}{n+2}+\cdots+\dfrac{1}{2n}\right).
\]
\end{example}
\exampleblank

\begin{example}
求
\[
    \lim\limits_{n\to\infty}
    \sqrt[n]{
        \prod\limits_{i=1}^{n}
        \dfrac{e^{1-\frac{1}{i}}}{\left(1+\dfrac{1}{i}\right)^i}
    }.
\]
\end{example}
\exampleblank

\begin{example}
求下列极限:
\begin{enumerate}[label=(\arabic*)]
    \item $\displaystyle
    \lim\limits_{x\to0}
    \dfrac{(1+x)^{\frac{1}{x}}-e}{x}$;
    \item $\displaystyle
    \lim\limits_{x\to0}
    \dfrac{(1+x)^{\frac{1}{x}}-(1+2x)^{\frac{1}{2x}}}{\sin x}$;
    \item $\displaystyle
    \lim\limits_{x\to+\infty}
    \left[
        \dfrac{x^2}{e}\left(1+\dfrac{1}{x}\right)^x
        -x^2+\dfrac{x}{2}
    \right]$.
\end{enumerate}
\end{example}
\exampleblank

\begin{example}
求
\[
    \lim\limits_{n\to\infty}\dfrac{(2n-1)!!}{(2n)!!}.
\]
\end{example}
\exampleblank

\newpage

\subsection{洛必达法则}

\begin{enumerate}[label=\textcircled{\arabic*},leftmargin=2.2em]
    \item 在使用洛必达法则之前,务必先检验是否属于可使用的几种不定式之一;
    \item 洛必达法则只是充分条件,而非必要条件,即就算洛必达法则算不出结果,也并不意味着极限不存在.例如
    $\displaystyle\lim\limits_{x\to+\infty}\dfrac{x+\sin x}{x+\cos x}=1$;
    \item 在使用洛必达法则前及过程中,能用等价无穷小代换的尽量先换,不要盲目一直求导.
\end{enumerate}

\begin{example}
求下列极限:
\begin{enumerate}[label=(\arabic*)]
    \item $\displaystyle
    \lim\limits_{x\to0}
    \left[\dfrac{1}{\ln(1+x)}-\dfrac{1}{x}\right]$;
    \item $\displaystyle
    \lim\limits_{x\to0}
    \left(\dfrac{\sin x}{x}\right)^{\frac{1}{1-\cos x}}$;
    \item $\displaystyle
    \lim\limits_{x\to\frac{\pi}{4}}
    \dfrac{\sec^2x-2\tan x}{1+\cos4x}$;
    \item $\displaystyle
    \lim\limits_{x\to+\infty}
    \dfrac{\left(\int_0^x e^{t^2}\,\mathrm{d}t\right)^2}
    {\int_0^x e^{2t^2}\,\mathrm{d}t}$;
    \item $\displaystyle
    \lim\limits_{x\to+\infty}\sqrt[x]{x}$.
\end{enumerate}
\end{example}
\exampleblank

\begin{example}
求下列极限:
\begin{enumerate}[label=(\arabic*)]
    \item $\displaystyle
    \lim\limits_{x\to+\infty}\left(\cos\dfrac{1}{x}+\sin\dfrac{1}{x}\right)$;
    \item $\displaystyle
    \lim\limits_{x\to0}
    \left[\dfrac{(1+x)^{\frac{1}{x}}}{e}\right]^{\frac{1}{x}}$;
    \item $\displaystyle
    \lim\limits_{x\to+\infty}
    \dfrac{\left(1+\dfrac{1}{x}\right)^x-e}{\sin\dfrac{1}{x}}$.
\end{enumerate}
\end{example}
\exampleblank

\begin{example}[\markstar]
已知函数 $f(x)$ 可导,且
\[
    \lim\limits_{x\to+\infty}\bigl[f(x)+f'(x)\bigr]=k,
\]
证明:
\[
    \lim\limits_{x\to+\infty}f(x)=k.
\]
\end{example}
\exampleblank

\newpage

\subsection{Taylor 公式}

事实上,Taylor 公式作为一元微分学巅峰,从证明过程来看,洛必达法则能求的极限,理论上 Taylor 公式都能求,但也仅限于理论上,因为实际计算中有时确实洛必达法则会更好算,这需根据具体情况而定.一个简单的判别方法(但不绝对):若所给函数求导后的结果相当复杂,就应优先 Taylor 公式,如复合了一堆函数.

\begin{example}
求下列极限:
\begin{enumerate}[label=(\arabic*)]
    \item $\displaystyle
    \lim\limits_{x\to0}
    \dfrac{\dfrac{x^2}{2}+1-\sqrt{1+x^2}}
    {(\cos x-e^{x^2})\sin x^2}$;
    \item $\displaystyle
    \lim\limits_{x\to0}
    \dfrac{e^x-1-x}{\sqrt{1-x}-\cos\sqrt{x}}$;
    \item $\displaystyle
    \lim\limits_{n\to\infty}
    \left(\dfrac{f\left(a+\dfrac{1}{n}\right)}{f(a)}\right)^n$,
    其中函数 $f(x)$ 在 $a$ 点可导,且 $f(a)\neq0$.
\end{enumerate}
\end{example}
\exampleblank

\begin{example}
求
\[
    \lim\limits_{n\to\infty}n\sin(2\pi n!e).
\]
\end{example}
\exampleblank

\begin{example}
求极限:
\begin{enumerate}[label=(\arabic*)]
    \item $\displaystyle
    \lim\limits_{x\to0}
    \dfrac{\tan(\tan x)-\sin(\sin x)}{x^3}$;
    \item $\displaystyle
    \lim\limits_{x\to0}
    \dfrac{\tan(\tan^2x)-x\sin(\sin x)}{x^4}$;
    \item $\displaystyle
    \lim\limits_{x\to0}
    \dfrac{x\sin(\sin x)-\sin^2x}{x^6}$.
\end{enumerate}
\end{example}
\exampleblank

\newpage

\subsection{迫敛性}

当所求函数极限比较复杂,不易直接求出时,可考虑适当的放缩,进而利用迫敛性.

\begin{example}
求下列极限:
\begin{enumerate}[label=(\arabic*)]
    \item $\displaystyle\lim\limits_{x\to0}x\left[\dfrac{1}{x}\right]$;
    \item $\displaystyle\lim\limits_{x\to+\infty}\dfrac{[x]}{x}$;
    \item $\displaystyle\lim\limits_{x\to+\infty}\dfrac{[xf(x)]}{x}$,
    其中 $\displaystyle\lim\limits_{x\to+\infty}f(x)=a$.
\end{enumerate}
\end{example}
\exampleblank

\begin{example}
已知 $a_i>0\ (i=1,2,\ldots,n)$,求
\[
    \lim\limits_{p\to+\infty}
    \left[
        \left(\sum\limits_{i=1}^n a_i^p\right)^{\frac{1}{p}}
        +
        \left(\sum\limits_{i=1}^n a_i^{-p}\right)^{\frac{1}{p}}
    \right].
\]
\end{example}
\exampleblank

\newpage

\subsection{Heine 归结原理}

利用 Heine 归结原理,可以沟通数列极限与函数极限,将两类问题进行转化.

\begin{example}
求下列极限:
\begin{enumerate}[label=(\arabic*)]
    \item $\displaystyle
    \lim\limits_{n\to\infty}
    \dfrac{(e-1)e^{\frac{1}{n}}}{n\left(e^{\frac{1}{n}}-1\right)}$;
    \item $\displaystyle
    \lim\limits_{n\to\infty}
    n\left(x^{\frac{1}{n}}-x^{\frac{1}{2n}}\right)\quad(x>0)$.
\end{enumerate}
\end{example}
\exampleblank

\begin{example}
求极限
\[
    \lim\limits_{n\to\infty}
    \sum\limits_{k=n}^{2n}\sin\dfrac{\pi}{k}.
\]
\end{example}
\exampleblank

\begin{example}
\begin{enumerate}[label=(\arabic*)]
    \item 设
    \[
        f(x)=\sum\limits_{k=1}^n a_k\sin kx
    \]
    满足 $|f(x)|\leq|\sin x|$,证明:
    \[
        \left|\sum\limits_{k=1}^n ka_k\right|\leq1.
    \]
    \item 设
    \[
        f(x)=\sum\limits_{k=1}^n a_k\ln(1+\ln x)
    \]
    满足 $|f(x)|\leq|x|$,且 $x>0$,证明:
    \[
        \left|\sum\limits_{k=1}^n ka_k\right|\leq1.
    \]
\end{enumerate}
\end{example}
\exampleblank

\newpage

\subsection{递推法}

\begin{example}
设 $f(x)$ 为 $(0,+\infty)$ 上函数,$a>0$ 且 $a\neq1$.
\begin{enumerate}[label=(\arabic*)]
    \item 若 $f(x)$ 满足 $f(x^a)=f(x)$,且
    \[
        \lim\limits_{x\to0^+}f(x)
        =\lim\limits_{x\to+\infty}f(x)
        =f(1),
    \]
    则 $f(x)\equiv f(1)$;
    \item 若 $f(x)$ 满足 $f(ax)=f(x)$,且
    $\displaystyle\lim\limits_{x\to0^+}f(x)$ 或
    $\displaystyle\lim\limits_{x\to+\infty}f(x)$ 存在,则 $f(x)$ 为常值函数.
\end{enumerate}
\end{example}
\exampleblank

\begin{example}
设 $f(x)$ 在 $\mathbb{R}$ 上有定义,在 $x=0$ 的某邻域有界,且
\[
    f(ax)=bf(x),
\]
其中 $a>1,b>1$,证明:
\[
    \lim\limits_{x\to0}f(x)=0.
\]
\end{example}
\exampleblank

\begin{example}
\begin{enumerate}[label=(\arabic*)]
    \item 设函数 $f:(0,+\infty)\to(0,+\infty)$ 单调递增,且
    \[
        \lim\limits_{t\to+\infty}\dfrac{f(2t)}{f(t)}=1,
    \]
    证明:对任意 $m>0$,有
    \[
        \lim\limits_{t\to+\infty}\dfrac{f(mt)}{f(t)}=1.
    \]
    \item 设 $\displaystyle\lim\limits_{x\to0}f(x)=0$,且
    \[
        \lim\limits_{x\to0}\dfrac{f(x)-f\left(\dfrac{x}{2}\right)}{x}=0,
    \]
    证明:
    \[
        \lim\limits_{x\to0}\dfrac{f(x)}{x}=0.
    \]
\end{enumerate}
\end{example}
\exampleblank

\newpage

\section{函数极限综合问题}

\newpage

\subsection{渐近线}

设 $f(x),g(x)$ 满足
$\displaystyle\lim\limits_{x\to x_0}[f(x)-g(x)]=0$,则也称 $g(x)$ 为 $f(x)$ 在 $x\to x_0$ 时的渐近曲线.一般情况下,常讨论 $g(x)$ 为一、二次曲线的情形.在一次曲线的情形,称 $g(x)$ 为渐近(直)线:$y=ax+b$.且有计算公式
\[
    a=\lim\limits_{x\to+\infty}\dfrac{f(x)}{x},
    \qquad
    b=\lim\limits_{x\to+\infty}[f(x)-ax].
\]
特别地:
\begin{enumerate}[label=\textcircled{\arabic*},leftmargin=2.2em]
    \item 若 $\displaystyle\lim\limits_{x\to+\infty}f(x)=a$ 或
    $\displaystyle\lim\limits_{x\to-\infty}f(x)=b$,则称 $y=a$ 或 $y=b$ 为 $y=f(x)$ 的水平渐近线;
    \item 若 $\displaystyle\lim\limits_{x\to x_0^+}f(x)=\pm\infty$ 或
    $\displaystyle\lim\limits_{x\to x_0^-}f(x)=\pm\infty$,则称 $x=x_0$ 为 $y=f(x)$ 的一条竖直渐近线.
\end{enumerate}

\begin{example}
求函数
\[
    y=\dfrac{(1+x)^2}{4(1-x)}
\]
的渐近线.
\end{example}
\exampleblank

\begin{example}
\begin{enumerate}[label=(\arabic*)]
    \item 求 $\alpha,\beta$ 值,使得
    \[
        \lim\limits_{x\to+\infty}
        \left(\sqrt{4x^2+x+1}-(\alpha x+\beta)\right)=0;
    \]
    \item 求 $\alpha,\beta,\gamma$ 表达式,使得
    \[
        \lim\limits_{x\to+\infty}
        [f(x)-(\alpha+\beta x+\gamma x^2)]=0.
    \]
\end{enumerate}
\end{example}
\exampleblank

\begin{example}
求曲线
\[
    (x-a)y^2=x^2(x-b),
    \qquad a>0,\ b>0,\ a\neq b,
\]
的渐近线.
\end{example}
\exampleblank

\newpage

\subsection{特殊函数性质}

\begin{proposition}[\markstar\ Riemann 函数]
设 Riemann 函数为
\[
R(x)=
\begin{cases}
    \dfrac{1}{q}, & x=\dfrac{p}{q},\\
    0, & x\text{ 为 }(0,1)\text{ 中的无理数或 }0,1,
\end{cases}
\]
则对任意 $x_0\in\lbrack 0,1\rbrack$,有
\[
    \lim\limits_{x\to x_0}R(x)=0.
\]
\end{proposition}

\begin{remark}
因此,Riemann 函数在每个点处极限值为零,进而其只在无理点连续(不考虑端点 $0,1$),有理点间断.
\end{remark}

\begin{example}
证明命题 2.2.1.
\end{example}
\exampleblank

\begin{proposition}[\markstar\ Dirichlet 函数]
设 Dirichlet 函数 $D:\mathbb{R}\to\{0,1\}$,
\[
D(x)=
\begin{cases}
    1, & x\in\mathbb{Q},\\
    0, & x\in\mathbb{R}\setminus\mathbb{Q}.
\end{cases}
\]
则对任意 $x_0\in\mathbb{R}$,
$\displaystyle\lim\limits_{x\to x_0}D(x)$ 不存在.
\end{proposition}

\begin{remark}
因此很自然可得 $D(x)$ 在任意一点不连续,由此还能构造只在某点连续,而其它点不连续的函数---$f(x)=(x-a)D(x)$.
\end{remark}

\begin{example}
证明命题 2.2.2.
\end{example}
\exampleblank

\begin{proposition}[\markstar\ 单调函数]
设 $f(x)$ 在 $\mathring{U}(x_0)$ 上为单调递增函数,则 $f(x_0-0),f(x_0+0)$ 均存在,且
\[
    f(x_0-0)=\sup\limits_{x\in\mathring{U}_{-}(x_0)}f(x),
    \qquad
    f(x_0+0)=\inf\limits_{x\in\mathring{U}_{+}(x_0)}f(x).
\]
\end{proposition}

\begin{remark}
因此单调函数一定存在单侧极限,进而间断点必为第一类间断点(跳跃间断点).
\end{remark}

\begin{example}
证明命题 2.2.3.
\end{example}
\exampleblank

\newpage

\subsection{推广的 Stolz 定理}

\begin{theorem}[\markstar\ 函数版本 Stolz]
设 $f(x),g(x)$ 在 $\lbrack a,+\infty)$ 上有定义,$T$ 是一个正常数,且满足:
\begin{enumerate}[label=(\arabic*)]
    \item $f(x),g(x)$ 在任一区间 $\lbrack a,b\rbrack$ 上有界;
    \item 对任意 $x>a$,有 $g(x+T)>g(x)$ 且 $g(x)\to+\infty\ (x\to+\infty)$.
\end{enumerate}
若存在极限
\[
    \lim\limits_{x\to+\infty}
    \dfrac{f(x+T)-f(x)}{g(x+T)-g(x)}=A,
\]
则有
\[
    \lim\limits_{x\to+\infty}\dfrac{f(x)}{g(x)}=A.
\]
\end{theorem}

\begin{example}
证明定理 2.2.1.
\end{example}
\exampleblank

\begin{example}[Cauchy 定理]
若 $f$ 在 $(a,+\infty)$ 内有定义,且内闭有界(即对任意 $\lbrack \alpha,\beta\rbrack\subseteq(a,+\infty)$,$f$ 在 $\lbrack \alpha,\beta\rbrack$ 上有界),则:
\begin{enumerate}[label=(\arabic*)]
    \item 当右边极限存在时,
    \[
        \lim\limits_{x\to+\infty}\dfrac{f(x)}{x}
        =\lim\limits_{x\to+\infty}[f(x+1)-f(x)];
    \]
    \item 当右边极限存在时,
    \[
        \lim\limits_{x\to+\infty}[f(x)]^{\frac{1}{x}}
        =\lim\limits_{x\to+\infty}\dfrac{f(x+1)}{f(x)},
        \qquad f(x)\geq c>0.
    \]
\end{enumerate}
\end{example}
\exampleblank

\begin{example}
设 $f(x)$ 在 $\lbrack a,+\infty)$ 上有定义,且内闭有界,
\[
    \lim\limits_{x\to+\infty}
    \dfrac{f(x+1)-f(x)}{x^n}=l
\]
($l$ 有限),证明:
\[
    \lim\limits_{x\to+\infty}\dfrac{f(x)}{x^{n+1}}
    =\dfrac{l}{n+1}.
\]
\end{example}
\exampleblank

\newpage

\subsection{周期函数}

\begin{example}
设 $f,g$ 是 $(-\infty,+\infty)$ 上的周期函数,若
\[
    \lim\limits_{x\to+\infty}[f(x)-g(x)]=0,
\]
则 $f(x)=g(x)$.
\end{example}
\exampleblank

\begin{example}
设 $f(x)$ 是周期为 $T\ (T>0)$ 的连续周期函数,则
\[
    \lim\limits_{x\to+\infty}
    \dfrac{1}{x}\int_0^x f(t)\,\mathrm{d}t
    =\dfrac{1}{T}\int_0^T f(t)\,\mathrm{d}t.
\]
\end{example}
\exampleblank

\newpage

\section{实数基本定理 *}

\begin{example}
设 $f$ 在有限区间 $I$ 上有定义,满足:对任意 $x\in I$,存在 $x$ 的邻域 $(x-\delta,x+\delta)$,使得 $f(x)$ 在 $(x-\delta,x+\delta)\cap I$ 上有界.
\begin{enumerate}[label=(\arabic*)]
    \item 证明:当 $I=\lbrack a,b\rbrack$ 时,$f(x)$ 在 $I$ 上有界;
    \item 当 $I=(a,b)$ 时,$f(x)$ 在 $I$ 上一定有界吗?
\end{enumerate}
\end{example}
\exampleblank

\begin{example}
设 $f(x)$ 在 $\lbrack a,b\rbrack$ 上有定义且在每一点处函数的极限存在,求证:$f(x)$ 在 $\lbrack a,b\rbrack$ 上有界.
\end{example}
\exampleblank

\begin{example}
若 $f(x)$ 在 $\lbrack a,b\rbrack$ 无界,试证明存在 $x_0\in\lbrack a,b\rbrack$,对任意 $\delta>0$,使得 $f(x)$ 在 $(x_0-\delta,x_0+\delta)$ 上无界.
\end{example}
\exampleblank

\begin{example}
设 $f(x)$ 在 $(a,b)$ 内有定义,对任意 $\xi\in(a,b)$,存在 $\delta>0$,当 $x\in(\xi-\delta,\xi+\delta)\cap(a,b)$ 时,有 $f(x)<f(\xi)$ 当 $x<\xi$ 时,$f(x)>f(\xi)$ 当 $x>\xi$ 时.求证:$f(x)$ 在 $(a,b)$ 内严格递增.
\end{example}
\exampleblank

\begin{example}
设 $f(x)$ 在 $\lbrack 0,1\rbrack$ 上递增,$f(a)>a,f(b)>b$,证明:
\begin{enumerate}[label=(\arabic*)]
    \item 存在 $x_0\in(a,b)$,使得 $f(x_0)=x_0$;
    \item 存在 $x_1\in(a,b)$,使得 $f(x_1)=x_1^2$.
\end{enumerate}
\end{example}
\exampleblank

\begin{example}
设定义在 $\lbrack 0,1\rbrack$ 上的 $f(x),g(x)$ 满足
\[
    f(0)>0,\qquad f(1)<1,\qquad g\in C(\lbrack 0,1\rbrack).
\]
若 $f+g$ 在 $\lbrack 0,1\rbrack$ 上递增,则存在 $\xi\in\lbrack 0,1\rbrack$,使得 $f(\xi)=0$.
\end{example}
\exampleblank

\begin{example}
设定义在 $(-\infty,+\infty)$ 上的 $f(x)$ 满足:对任意 $x_0\in(-\infty,+\infty)$,存在 $\delta>0$,使得
\[
    f(x_0)\geq f(x),
    \qquad x\in(x_0-\delta,x_0+\delta).
\]
证明:存在区间 $I$,使得 $f(x)$ 在 $I$ 上是常数.
\end{example}
\exampleblank
